\documentclass[11pt]{article}

\usepackage[utf8]{inputenc}
\usepackage[T1]{fontenc}
\usepackage{lmodern}
\usepackage{amsmath,amssymb,amsthm,mathtools}
\usepackage{stmaryrd}
\usepackage[margin=1in]{geometry}
\usepackage{enumitem}
\usepackage{xcolor}
\usepackage{hyperref}
\hypersetup{
  colorlinks=true,
  linkcolor=blue!60!black,
  citecolor=blue!60!black,
  urlcolor=blue!60!black
}

% theorem environments
\newtheorem{theorem}{Theorem}
\newtheorem{lemma}[theorem]{Lemma}
\newtheorem{proposition}[theorem]{Proposition}
\newtheorem{corollary}[theorem]{Corollary}
\newtheorem{question}[theorem]{Question}
\newtheorem{conjecture}[theorem]{Conjecture}
\theoremstyle{definition}
\newtheorem{definition}[theorem]{Definition}
\newtheorem{example}[theorem]{Example}
\newtheorem{remark}[theorem]{Remark}

\usepackage[nameinlink,capitalize]{cleveref}
\crefname{theorem}{Theorem}{Theorems}
\crefname{lemma}{Lemma}{Lemmas}
\crefname{proposition}{Proposition}{Propositions}
\crefname{corollary}{Corollary}{Corollaries}
\crefname{definition}{Definition}{Definitions}
\crefname{remark}{Remark}{Remarks}
\Crefname{theorem}{Theorem}{Theorems}
\Crefname{lemma}{Lemma}{Lemmas}
\Crefname{proposition}{Proposition}{Propositions}
\Crefname{corollary}{Corollary}{Corollaries}
\Crefname{definition}{Definition}{Definitions}
\Crefname{remark}{Remark}{Remarks}

\sloppy
\setlength{\emergencystretch}{3em}

% macros
\newcommand{\R}{\mathbb{R}}
\newcommand{\N}{\mathbb{N}}
\newcommand{\E}{\mathbb{E}}
\newcommand{\Prob}{\mathbb{P}}
\newcommand{\A}{\mathcal{A}}
\newcommand{\F}{\mathcal{F}}
\newcommand{\K}{\mathcal{K}}
\newcommand{\T}{\mathcal{T}}
\newcommand{\Q}{\mathcal{Q}}
\newcommand{\Sset}{\mathcal{S}}
\newcommand{\U}{\mathcal{U}}
\newcommand{\X}{\mathcal{X}}
\newcommand{\C}{\mathcal{C}}
\newcommand{\Homp}{\operatorname{Hom}^{+}}
\newcommand{\Homalg}{\operatorname{Hom}}
\newcommand{\Ext}{\operatorname{Ext}}
\newcommand{\supp}{\operatorname{supp}}
\newcommand{\Forb}{\operatorname{Forb}}
\newcommand{\brak}[1]{\left\langle #1\right\rangle}
\newcommand{\avg}[2]{\left\llbracket #1 \right\rrbracket_{#2}}
\newcommand{\qmap}{\mathfrak q}
\newcommand{\eps}{\varepsilon}
\newcommand{\one}{\mathbf 1}
\DeclareMathOperator{\cl}{cl}

\title{Root Planting, Quotient Semantics, and a Counterexample\\
       for Forbidden-Subgraph Reasoning in Flag Algebras}
\author{}
\date{\today}

\begin{document}
\maketitle

\begin{abstract}
Forbidden-subgraph reasoning in Razborov's flag algebras can be carried out in two
ways: in the \emph{quotient} flag algebra of the constrained class, the algebra one
actually computes in, or by \emph{ensemble} semantics in the ambient class of all graphs, which
asks that an inequality hold almost surely along the random labellings of unlabelled
constrained limits.  Quotient
non-negativity always implies ensemble non-negativity, so the quotient calculus is
unconditionally sound; we study the converse, its completeness.  We prove that the two
semantics agree, for every flag-algebra element, exactly when a support-closure
condition, which we call \emph{root-plantability}, holds: every labelled constrained limit
must be approximable by the labelled views of unlabelled constrained limits. Closure under
independent blow-ups (which holds for the clique-free classes, including triangle-free) implies root-plantability and ensures that these two notions of non-negativity coincide. But heredity alone does not lead to root-plantability.  Every class so
sparse that its $n$-vertex graphs have only $o(n^2)$ edges, yet that still contains
arbitrarily large stars---the $C_4$-free, $K_{s,t}$-free, $C_{2k}$-free, and planar
graphs among them---is not root-plantable at the one-vertex type; by complementation neither are their dense
complements, so neither sparsity nor density is the dividing line.  The uniform cause is
\emph{pinning}: a labelled quantity that takes one fixed value across all constrained
limits while the constrained class admits a labelled limit where it takes another.  We conjecture
that pinning is the basic obstruction, but the $C_5$-free graphs show that it must be
analysed type by type.  The one-vertex and two-root non-edge types admit local planting
constructions; at the two-root edge type a book edge gives a new pinning obstruction, so
the all-types $C_5$-free conjecture is false.

This incompleteness, however, never affects a density bound of the kind a flag-algebra
computation is run to produce.  Such a bound is an inequality about unlabelled graphs, whereas
the gap between the two semantics lives entirely in the auxiliary inequalities---about
vertex-labelled graphs---that a proof passes through on the way to it.  We show
that on unlabelled-graph statements the two semantics always coincide, that the labelled
examples witnessing the gap vanish once averaged back into unlabelled-graph statements, and
that the gap can open only at a degenerate extreme.  So for every class we exhibit, the
incompleteness leaves the final density bounds untouched, and whether any class lets it
reach one remains open.
\end{abstract}

\section{Introduction}

A recurring task in extremal combinatorics is to prove a density inequality that
holds only under a forbidden-subgraph constraint, such as ``in every triangle-free
graph this combination of subgraph densities is non-negative.''  Razborov's flag
algebras offer two routes to such statements.  One may pass to the \emph{quotient}
flag algebra of the constrained class and ask for non-negativity there
(\emph{quotient semantics}); or one may stay in the ambient class of all graphs and
ask that the inequality hold almost surely along the random labelled extensions of
constrained limits (\emph{ensemble semantics}).  The first is the algebra one
actually computes in; the second is what we want to hold throughout the
constrained class.  When the two agree, a certificate proved in the quotient is a
genuine density bound for the constrained class.  This paper determines exactly when
they agree.

One direction is automatic.  Quotient semantics always implies ensemble semantics
(\cref{thm:support-criterion}), so a quotient certificate is unconditionally
\emph{sound}: the validity of forbidden-subgraph flag-algebra reasoning is never in
question.  The substantive question is the converse---whether \emph{every} true
constrained density inequality can be certified in the quotient---so our subject is the
\emph{completeness} of the quotient calculus.

Let $\T_0$ be the class of all finite simple graphs, and let $\T_1$ be a
\emph{hereditary} subclass: a class of finite graphs closed under taking induced
subgraphs, equivalently the graphs that avoid every member of some forbidden family.
The reader should keep in mind a forbidden-subgraph class such as the triangle-free or
$C_4$-free graphs.  For
each non-degenerate type $\sigma$ of $\T_1$, Razborov's construction gives a
quotient homomorphism
\[
  \qmap_\sigma:\A^\sigma[\T_0]\twoheadrightarrow \A^\sigma[\T_1].
\]
The quotient semantics of an element $f\in\A^\sigma[\T_0]$ says that
\[
  \psi(\qmap_\sigma f)\ge 0
  \qquad\text{for every }\psi\in\Homp(\A^\sigma[\T_1],\R).
\]
The ensemble semantics stays in the ambient class $\T_0$ of all graphs.  It asks that, for
every unlabelled $\T_1$-supported homomorphism $\phi_0$, the canonical random
extension $\phi^\sigma$ of $\phi_0$ to type $\sigma$ satisfies
\[
  \phi^\sigma(f)\ge 0
  \qquad\text{almost surely.}
\]

\paragraph{The criterion.}
The right property is not simply heredity, but the following support
condition:
\[
  \overline{\bigcup_{\substack{\phi_0\ \T_1\text{-supported,}\\ \text{positive }\sigma\text{-density}}}\;\supp\bigl(\Ext_\sigma(\phi_0)\bigr)}
  \;=\;
  \{\psi\circ\qmap_\sigma:\psi\in\Homp(\A^\sigma[\T_1],\R)\}.
\]
The union ranges only over the unlabelled $\T_1$-supported homomorphisms $\phi_0$ of
positive $\sigma$-density---those for which the canonical random extension above is
defined; $\Ext_\sigma(\phi_0)$ denotes the law of that extension, $\supp$ its
(topological) support, and the closure is taken in the compact space
$\Homp(\A^\sigma[\T_0],\R)$ under pointwise convergence of flag densities.  Both sides lie
in that ambient space: each $\psi\circ\qmap_\sigma$ is again a positive homomorphism on
$\A^\sigma[\T_0]$, since $\qmap_\sigma$ carries every $\T_0$-flag to a $\T_1$-flag or to
$0$, so the composite stays non-negative on flags.  (\Cref{sec:background} and
\cref{def:root-planting} make all of this precise.)  Informally, every constrained
labelled limit must be approximable by positive-probability labelled views of constrained
unlabelled limits.  We call this condition \emph{root-plantability} and prove that it is
exactly equivalent to the quotient/ensemble correspondence for every flag-algebra element
(\cref{thm:support-criterion}).

\paragraph{Results.}
Root-plantability is implied by closure under independent blow-ups, which recovers the
correspondence for the clique-free classes, including triangle-free.  Heredity alone, however,
does not imply it: every \emph{edge-degenerate} class---one whose $n$-vertex graphs have
only $o(n^2)$ edges---that contains arbitrarily large stars (the $C_4$-free,
$K_{s,t}$-free, $C_{2k}$-free, and planar graphs among them) fails at the one-vertex
type, witnessed by the labelled edge flag, and by complementation so do their dense
complements; neither sparsity nor density is the dividing line.  The uniform cause is
\emph{pinning}: a labelled quantity that takes one fixed value across all constrained
limits while the constrained class admits a labelled limit where it takes another.  We conjecture
that pinning characterises root-plantability, understood type by type.  The $C_5$-free
graphs now illustrate both sides: the one-vertex and two-root non-edge types can be
planted by local repairs of the blow-up construction, while the two-root edge type has a
book-edge pinning obstruction.  This incompleteness, moreover, never
reaches a density bound.  A density bound is a statement about unlabelled graphs---the
\emph{empty type}---where the two semantics provably coincide (\cref{prop:empty-type});
the labelled examples that witness the gap vanish when averaged back to unlabelled graphs
(\cref{prop:ideal-zero}); and the gap can open only at a degenerate extreme, since for
subgraph-forbidden classes pinning is always to a boundary value (\cref{thm:no-interior}).  The incompleteness is
thus confined to the intermediate, labelled certificates a proof passes through; for the
obstructions we exhibit it never affects a density bound (\cref{sec:empty-type}), and
whether some constrained class with a sufficiently rich labelled support makes it do so is
left open.

\section{Flag-algebra background}\label{sec:background}

We recall the flag-algebra apparatus \cite{Razborov2007}, at the level of detail a reader
new to it needs.  Everything is phrased for finite simple graphs.  Throughout, $\T$ stands
for a hereditary class of finite simple graphs---either the class $\T_0$ of all finite
graphs, or a forbidden-subgraph subclass $\T_1$---and ``graph in $\T$'' means a finite
graph belonging to that class.

\paragraph{Types and flags.}
A \emph{type} $\sigma$ is a fixed graph in $\T$ on $k$ \emph{labelled} vertices
(its \emph{roots})---$k$ distinguished vertices together with the edges among
them.  A \emph{$\sigma$-flag} is a graph in $\T$ with a distinguished labelled copy
of $\sigma$ singled out inside it---informally, a finite graph $k$ of whose vertices
are marked so as to form the pattern $\sigma$.  The \emph{empty type} ($k=0$) marks
nothing, so a flag of the empty type is just a graph in $\T$.  Flags are taken up to
label-preserving isomorphism, and $\F^\sigma[\T]$ (abbreviated $\F^\sigma$) is the
countable set of all of them.  For $\sigma$-flags $F$ and $G$ with $|F|\le|G|$, the
\emph{density} $p(F,G)\in[0,1]$ is the probability that a uniformly random set of $|F|-k$
unlabelled vertices of $G$, taken together with the $k$ roots, induces a copy of $F$; for
the empty type this is the ordinary induced-subgraph density.

\paragraph{The flag algebra.}
$\A^\sigma[\T]$ is the real vector space spanned by $\F^\sigma$ (formal real combinations
of flags), modulo the linear relations
\[
  F=\sum_{|H|=|F|+1}p(F,H)\,H
\]
relating each flag to the flags one vertex larger.  These relations encode the
consistency of densities across sizes---$p(F,G)=\sum_{|H|=|F|+1}p(F,H)\,p(H,G)$ for every
larger $G$---and are what make ``density'' a well-defined linear functional on the whole
space.  The space
carries a commutative, associative \emph{product} that records how densities combine: the
product $F_1\cdot F_2$ of $\sigma$-flags on $n_1$ and $n_2$ vertices is the formal
combination of the $\sigma$-flags $H$ on $n_1+n_2-k$ vertices, each weighted by the
probability that splitting the unlabelled vertices of $H$ uniformly into parts of sizes
$n_1-k$ and $n_2-k$ exhibits $F_1$ on one part (with the roots) and $F_2$ on the other.
The product is rigged so that densities multiply in the limit (made precise next); its
unit is the flag $\one_\sigma$ consisting of the type alone.  We write $\A^0[\T]$ for the
empty-type algebra.

\paragraph{Positive homomorphisms as limits.}
A \emph{positive homomorphism} is a linear, unit-preserving, multiplicative map
$\phi:\A^\sigma[\T]\to\R$---so $\phi(\one_\sigma)=1$ and
$\phi(F_1\cdot F_2)=\phi(F_1)\,\phi(F_2)$---that is non-negative on flags, i.e.\
$\phi(F)\ge0$ for every $F\in\F^\sigma$.  The set of all of them is denoted
$\Homp(\A^\sigma[\T],\R)$.  By Razborov's theory these are exactly the limits of
convergent flag sequences: a sequence of finite $\sigma$-flags of growing size
\emph{converges} when every density $p(F,\cdot)$ converges, its limit is the homomorphism
$\phi$ with $\phi(F)=\lim p(F,\cdot)$, and conversely every positive homomorphism arises
this way \cite[Theorem 3.3]{Razborov2007}.  We accordingly picture a positive homomorphism
$\phi$ as a limit of finite graphs and $\phi(F)$ as the density of the flag $F$ in it, and
we use the terms \emph{positive homomorphism} and \emph{limit} interchangeably.  On a single
flag this value lies in $[0,1]$: for each size the flags of that size sum to the unit,
$\sum_{|H|=|F|}H=\one_\sigma$ (a vertex set induces exactly one flag), so
$0\le\phi(F)\le\sum_{|H|=|F|}\phi(H)=\phi(\one_\sigma)=1$.  Hence evaluation on all flags,
$\phi\mapsto(\phi(F))_{F\in\F^\sigma}$, embeds $\Homp(\A^\sigma[\T],\R)$ into the compact
metrizable cube $[0,1]^{\F^\sigma[\T]}$; the image is closed---cut out by the linear and
multiplicative homomorphism conditions---so $\Homp(\A^\sigma[\T],\R)$ is itself compact
metrizable.

The averaging (unlabelling) operator
\[
  \avg{\cdot}{\sigma}:\A^\sigma[\T]\to\A^0[\T]
\]
forgets the labels of a $\sigma$-flag and averages over the choice of labelling.  The
only properties we use are that it is linear and that on a single flag $F$ with
underlying unlabelled graph $M$ it returns a \emph{non-negative scalar multiple}
of $M$.  Recall that $\one_\sigma\in\A^\sigma[\T]$, the flag consisting of the type alone, is the
multiplicative unit of $\A^\sigma[\T]$.  Its unlabelling
\[
  \brak{\sigma}:=\avg{\one_\sigma}{\sigma}\in\A^0[\T]
\]
is the unlabelled density of embeddings of the type $\sigma$: for a limit
$\phi_0$, the value $\phi_0(\brak\sigma)$ is the probability that $k$ random vertices, in
a random order, form a copy of $\sigma$.  The type $\sigma$ is
\emph{non-degenerate} if $\brak\sigma\neq0$ in $\A^0[\T]$.  By the unlabelling property
above, $\brak\sigma=\avg{\one_\sigma}{\sigma}$ is a non-negative multiple of the unlabelled
graph underlying $\sigma$, so it is a \emph{non-negative element}: $\phi_0(\brak\sigma)\ge0$
for every limit $\phi_0$.  Since an
element of $\A^0[\T]$ is the zero element exactly when every limit assigns it the value
$0$, $\brak\sigma\neq0$ holds precisely when $\phi_0(\brak\sigma)>0$ for some limit
$\phi_0$.  If
$\phi_0\in\Homp(\A^0[\T],\R)$ and $\phi_0(\brak\sigma)>0$, Razborov's extension
theorem gives a unique probability measure on
$\Homp(\A^\sigma[\T],\R)$, denoted here by
\[
  \Ext_\sigma(\phi_0),
\]
characterised by
\begin{equation}\label{eq:extension-expectation}
  \E_{\phi^\sigma\sim\Ext_\sigma(\phi_0)}[\phi^\sigma(g)]
  =
  \frac{\phi_0(\avg{g}{\sigma})}{\phi_0(\brak\sigma)}
  \qquad(g\in\A^\sigma[\T]).
\end{equation}
Intuitively, $\Ext_\sigma(\phi_0)$ is the law of the labelled limit seen from a uniformly
random copy of $\sigma$ inside $\phi_0$; \eqref{eq:extension-expectation} records that
the expected labelled density of $g$ is the unlabelled density of $g$, normalised by that
of $\sigma$.  Moreover, if a finite sequence converges to $\phi_0$, then the
corresponding finite random-labelling distributions converge weakly to $\Ext_\sigma(\phi_0)$
\cite[Theorems 3.5 and 3.12]{Razborov2007}.

Recall that the \emph{support} of a Borel probability measure $P$ on a compact metrizable
space $X$, written $\supp(P)$, is the smallest closed set of full $P$-measure; equivalently,
\[
  \supp(P)=\{x\in X:P(U)>0\text{ for every open }U\ni x\}.
\]
We will repeatedly use the following elementary consequence.

\begin{lemma}[Almost-sure non-negativity and support]\label{lem:support-as}
Let $P$ be a Borel probability measure on a compact metrizable space $X$, and let
$h:X\to\R$ be continuous.  Then
\[
  P[h\ge0]=1
\]
if and only if $h(x)\ge0$ for every $x\in\supp(P)$.
\end{lemma}

\begin{proof}
($\Rightarrow$) Suppose $P[h\ge0]=1$, equivalently $P[h<0]=0$, and let
$x\in\supp(P)$.  If $h(x)<0$, then $\{y:h(y)<0\}$ is an open set (as $h$ is continuous)
containing $x$, so by the defining property of the support it has positive $P$-measure,
contradicting $P[h<0]=0$.  Hence $h(x)\ge0$.

($\Leftarrow$) Suppose $h\ge0$ on $\supp(P)$.  Then
$\{y:h(y)\ge0\}\supseteq\supp(P)$, and the support has full measure, so
$P[h\ge0]\ge P(\supp(P))=1$; that is, $P[h\ge0]=1$.
\end{proof}

\section{Constrained classes and supported homomorphisms}

Let $\T_1$ be a hereditary subclass of $\T_0$.  For a non-degenerate type $\sigma$ of
$\T_1$, define
\[
  X_\sigma:=\Homp(\A^\sigma[\T_0],\R).
\]
Call a $\sigma$-flag \emph{forbidden} when its underlying unlabelled graph belongs to
$\T_0$ but not to $\T_1$.  Being forbidden depends only on
that underlying graph, not on how the flag is labelled.  The quotient map
\[
  \qmap_\sigma:\A^\sigma[\T_0]\twoheadrightarrow\A^\sigma[\T_1]
\]
is precisely the map that sends each forbidden flag to $0$ and every other flag to itself.
It pulls each $\psi\in\Homp(\A^\sigma[\T_1],\R)$ back to a point of $X_\sigma$, the composite
${\psi} \circ {\qmap_\sigma}$: this is an algebra homomorphism, and it is non-negative on every
flag because $\qmap_\sigma$ sends flags to flags or to $0$ and $\psi$ is non-negative on
flags.  We write
\[
  Q_\sigma
  :=
  \{\psi\circ\qmap_\sigma:
      \psi\in\Homp(\A^\sigma[\T_1],\R)\}
  \subseteq X_\sigma.
\]
Thus $Q_\sigma$ is the constrained, or $\T_1$-supported, part of the ambient homomorphism
space.  It has a simple intrinsic description, phrased entirely in terms of $\chi$: a
homomorphism $\chi\in X_\sigma$ lies in $Q_\sigma$ if and only if it vanishes on every
forbidden flag.  One direction is immediate.  Since $\qmap_\sigma$ sends each forbidden
flag to $0$, any $\chi=\psi\circ\qmap_\sigma$ satisfies
$\chi(F)=\psi(\qmap_\sigma(F))=\psi(0)=0$ on every forbidden flag $F$, so every member of
$Q_\sigma$ vanishes on the forbidden flags.

The converse---that a homomorphism vanishing on all forbidden flags is itself of the form
$\psi\circ\qmap_\sigma$---is the substantive direction, and it is where heredity does the
work.  Forming $\A^\sigma[\T_1]$ amounts to declaring the forbidden flags equal to $0$, and
heredity makes this consistent with multiplication: the product of a forbidden flag with
any flag is again a combination of forbidden flags, because any graph that contains a
non-$\T_1$ graph as an induced subgraph is itself outside $\T_1$.  Hence the elements that
$\qmap_\sigma$ collapses to $0$ are exactly the linear combinations of forbidden flags.  A
homomorphism $\chi$ that vanishes on all of them cannot distinguish $\A^\sigma[\T_0]$ from
its quotient, so it descends through $\qmap_\sigma$: there is a unique homomorphism $\psi$
on $\A^\sigma[\T_1]$ with $\chi=\psi\circ\qmap_\sigma$, and $\psi$ is non-negative on flags
because $\chi$ is.  Thus $\psi\in\Homp(\A^\sigma[\T_1],\R)$ and $\chi\in Q_\sigma$.

This description also makes the closedness of $Q_\sigma$ transparent.  For a single
forbidden flag $F$, the condition $\chi(F)=0$ is a \emph{closed} one: the evaluation
$\chi\mapsto\chi(F)$ is continuous---recall that $X_\sigma$ carries the topology of
pointwise convergence of flag densities (\cref{sec:background})---so the homomorphisms
satisfying it form the preimage $\{\chi:\chi(F)=0\}$ of the closed set $\{0\}$.  The space
$Q_\sigma$ consists of the homomorphisms that satisfy all of these conditions at once, one
per forbidden flag, so it is an intersection of closed sets and is therefore closed in
$X_\sigma$.

Informally, the next lemma says that randomly labelling a constrained unlabelled limit
almost surely yields a constrained labelled one.

\begin{lemma}[Support passes to random extensions]\label{lem:support-passes-general}
Let $\phi_0\in Q_0$ and suppose $\phi_0(\brak\sigma)>0$.  If
$\phi^\sigma\sim\Ext_\sigma(\phi_0)$ is formed in the ambient class $\T_0$,
then
\[
  \phi^\sigma\in Q_\sigma
  \qquad\text{almost surely.}
\]
\end{lemma}

\begin{proof}
Let $F$ be a forbidden $\sigma$-flag, and write $M\notin\T_1$ for its underlying unlabelled
graph.  The downward average $\avg{F}{\sigma}$ is a non-negative scalar multiple of $M$ as
an element of $\A^0[\T_0]$.  Since $\phi_0\in Q_0$ vanishes on every graph outside $\T_1$,
we have $\phi_0(M)=0$, and hence $\phi_0(\avg{F}{\sigma})=0$.  By \eqref{eq:extension-expectation},
\[
  \E[\phi^\sigma(F)]=0.
\]
But $\phi^\sigma(F)\ge0$ pointwise.  Hence $\phi^\sigma(F)=0$ almost surely.
Intersecting the resulting probability-one events over the countable set of
forbidden $\sigma$-flags shows that $\phi^\sigma\in Q_\sigma$ almost surely.
\end{proof}

\section{The support-closure criterion}

The following definition is the main conceptual point of this paper.

\begin{definition}[Root-planting set and root-plantability]\label{def:root-planting}
Fix a non-degenerate type $\sigma$ of $\T_1$.  Define
\[
  S_\sigma
  :=
  \overline{
  \bigcup_{\substack{\phi_0\in Q_0\\ \phi_0(\brak\sigma)>0}}
  \supp(\Ext_\sigma(\phi_0))}
  \subseteq X_\sigma,
\]
where the closure is taken in the compact ambient space
$X_\sigma=\Homp(\A^\sigma[\T_0],\R)$.
We say that the triple $(\T_0,\T_1,\sigma)$ is \emph{root-plantable} if
\[
  S_\sigma=Q_\sigma.
\]
\end{definition}

By \cref{lem:support-passes-general}, one always has $S_\sigma\subseteq Q_\sigma$.
The point of root-plantability is the reverse inclusion: every constrained
labelled homomorphism should be approximable by support points of random
extensions from constrained unlabelled homomorphisms.

\begin{theorem}[Support-closure criterion]\label{thm:support-criterion}
Let $\T_1$ be a hereditary subclass of $\T_0$, and let $\sigma$ be a non-degenerate type of $\T_1$.  For $f\in\A^\sigma[\T_0]$, consider the following two
conditions.
\begin{enumerate}[label=\textup{(\alph*)}, leftmargin=2.8em]
  \item\label{cond:quotient}
  Quotient semantics:
  \[
    \chi(f)\ge0\qquad\text{for every }\chi\in Q_\sigma.
  \]
  Equivalently, $\psi(\qmap_\sigma f)\ge0$ for every
  $\psi\in\Homp(\A^\sigma[\T_1],\R)$.
  \item\label{cond:ensemble}
  Ensemble semantics:
  for every $\phi_0\in Q_0$ with $\phi_0(\brak\sigma)>0$,
  \[
    \Prob_{\phi^\sigma\sim\Ext_\sigma(\phi_0)}[\phi^\sigma(f)\ge0]=1.
  \]
\end{enumerate}
Then \ref{cond:quotient} always implies \ref{cond:ensemble}.  Moreover,
\ref{cond:quotient} and \ref{cond:ensemble} are equivalent for every
$f\in\A^\sigma[\T_0]$ if and only if $(\T_0,\T_1,\sigma)$ is root-plantable.
\end{theorem}

In words: a certificate proved in the quotient is always a valid constrained density
bound, and it captures \emph{every} such bound precisely when $\sigma$ is
root-plantable; when root-plantability fails, some inequality that holds almost surely
under random labelling cannot be seen in the quotient.

\begin{proof}
By \cref{lem:support-passes-general}, every support appearing in the definition
of $S_\sigma$ is contained in $Q_\sigma$.  Hence non-negativity on $Q_\sigma$
implies non-negativity on all those supports, and by \cref{lem:support-as} this
is precisely the implication from quotient semantics to ensemble semantics.

For the converse direction, note first that ensemble semantics for a fixed $f$
is equivalent, again by \cref{lem:support-as}, to non-negativity of $f$ on the
union of the supports in \cref{def:root-planting}.  Since $f$ is continuous on
$X_\sigma$, this is equivalent to non-negativity on $S_\sigma$.
Thus the comparison is simply
\[
  f\ge0\text{ on }S_\sigma
  \qquad\text{versus}\qquad
  f\ge0\text{ on }Q_\sigma.
\]
If $S_\sigma=Q_\sigma$, the two conditions are equivalent for every $f$.

Conversely suppose $S_\sigma\ne Q_\sigma$.  Choose
$\psi\in Q_\sigma\setminus S_\sigma$.  The compact space $X_\sigma$ is a closed subspace of the product
$[0,1]^{\F^\sigma}$, hence is compact metrizable and normal.  The closed sets
$S_\sigma$ and $\{\psi\}$ are disjoint subsets of the closed space $Q_\sigma$.
By Urysohn's lemma there is a continuous function $h\in C(Q_\sigma)$ such that
$h=1$ on $S_\sigma$ and $h(\psi)=-1$.  By Tietze's extension theorem, extend
$h$ to a continuous function $H\in C(X_\sigma)$.  By the Stone-Weierstrass
theorem in the form used by Razborov \cite[Proposition 3.7]{Razborov2007},
$\A^\sigma[\T_0]$ is uniformly dense in $C(X_\sigma)$.  Choose
$f\in\A^\sigma[\T_0]$ with $\|f-H\|_\infty<1/3$.  Then $f>0$ on
$S_\sigma$, so ensemble semantics holds for $f$, while $f(\psi)<0$, so quotient
semantics fails.  Thus equivalence for all $f$ forces $S_\sigma=Q_\sigma$.
\end{proof}

\section{Clone-closed graph classes are root-plantable}

We now give a checkable sufficient condition for root-plantability of graph
classes, built on an independent blow-up construction.  The idea is that when a class is
closed under blow-ups, a labelled vertex can be duplicated into a clone class of
positive density, so a configuration the quotient exhibits only at an exceptional vertex
can instead be realised at a random---hence typical---vertex, which is just what
root-plantability requires.  Throughout, $\T_0=\T_{\mathrm{Graph}}$ denotes the class of
all finite simple graphs, and $\K$ is a hereditary class of finite graphs, identified with
the constrained class $\T_1=\T_\K$ of graphs avoiding every induced subgraph not in $\K$.

\begin{definition}[Independent blow-up]
Let $G$ be a finite graph and let $m_v\ge1$ be an integer for each
$v\in V(G)$.  The independent blow-up $G^{\mathbf m}$ is obtained by replacing
each vertex $v$ by an independent clone class $C_v$ of size $m_v$, and by putting
all edges between $C_v$ and $C_w$ whenever $vw\in E(G)$.
\end{definition}

\begin{definition}[Clone-closed hereditary class]
A hereditary graph class $\K$ is \emph{clone-closed} if every independent
blow-up of every graph in $\K$ again belongs to $\K$.
\end{definition}

\begin{example}
For every $r\ge3$, the class of $K_r$-free graphs is clone-closed.  A clique in
an independent blow-up uses at most one vertex from each clone class, and hence
projects to a clique of the same size in the original graph.  In particular, the
triangle-free class is clone-closed.
\end{example}

\subsection{The planted blow-up estimate}

Let $F=(G,\theta)$ be a $\sigma$-flag with $|\sigma|=k$ and $|G|=n>k$.  Fix a
parameter $0<\lambda\le 1/2$.  If $k>0$, assign weights
\[
  w_{\theta(i)}=\frac{\lambda}{k}\quad(1\le i\le k),
  \qquad
  w_v=\frac{1-\lambda}{n-k}\quad(v\notin\operatorname{im}\theta).
\]
If $k=0$, set $w_v=1/n$ for all $v\in V(G)$.  For a large integer $N$, let
$B_N(G,\theta,\lambda)$ be an independent blow-up of $G$ whose clone class
$C_v$ has size approximately $w_vN$, and put
\[
  \operatorname{err}_N
  :=
  \sum_{v\in V(G)}\left|\frac{|C_v|}{N}-w_v\right|.
\]
The error tends to zero along suitable choices of integer cluster sizes.

\begin{lemma}[Planted blow-up approximation]\label{lem:planted-estimate}
Let $m\ge k$ be fixed.  There is a constant $C_m$, depending only on $m$, such
that the following holds for every $0<\lambda\le 1/2$.  Let $F_0$ be any $\sigma$-flag with at most $m$
vertices, let $(G,\theta)$ be a $\sigma$-flag on $n>k$ vertices, and let
$B_N=B_N(G,\theta,\lambda)$.  If
$\widehat\theta:\sigma\to B_N$ satisfies
$\widehat\theta(i)\in C_{\theta(i)}$ for all $i$, then
\begin{equation}\label{eq:planted-estimate}
  \left|p(F_0,(B_N,\widehat\theta))-p(F_0,(G,\theta))\right|
  \le
  C_m\left(\lambda+\frac1{n-k}+\operatorname{err}_N\right).
\end{equation}
Consequently, for every fixed $f\in\A^\sigma[\T_{\mathrm{Graph}}]$ there is a
constant $C_f$ such that
\begin{equation}\label{eq:planted-linear-estimate}
  \left|p(f,(B_N,\widehat\theta))-p(f,(G,\theta))\right|
  \le
  C_f\left(\lambda+\frac1{n-k}+\operatorname{err}_N\right).
\end{equation}
\end{lemma}

\begin{proof}
It suffices to prove \eqref{eq:planted-estimate}; the estimate for $f$ follows
by linearity because $f$ is a finite linear combination of flags.

Let $\ell=|F_0|\le m$.  To compute $p(F_0,(B_N,\widehat\theta))$, choose
$\ell-k$ additional unlabelled vertices uniformly from $B_N$, avoiding the
labelled vertices.  Call the sample good if none of these additional vertices
lies in a planted labelled clone class $C_{\theta(1)},\ldots,C_{\theta(k)}$ and
no two of them lie in the same non-labelled clone class.  The probability of
hitting a planted labelled clone class is at most
$(\ell-k)\lambda+O_m(\operatorname{err}_N)$.  Conditional on avoiding the planted clone classes, the normalised target weights
of the non-labelled clone classes are equal, up to $O_m(\operatorname{err}_N)$,
because the total non-labelled mass is $1-\lambda\ge1/2$.  Therefore the
probability that two sampled vertices
fall in the same non-labelled clone class is at most
$\binom{\ell-k}{2}/(n-k)+O_m(\operatorname{err}_N)$.  Hence the bad probability is
bounded by the right-hand side of \eqref{eq:planted-estimate}, after enlarging
$C_m$.

On the good event, projecting every sampled vertex of $B_N$ to its base vertex
in $G$ gives a uniformly chosen $(\ell-k)$-subset of
$V(G)\setminus\operatorname{im}\theta$, up to total variation error
$O_m(\operatorname{err}_N)$.  The induced labelled graph in the blow-up is then
exactly the induced labelled graph in $(G,\theta)$ on the projected vertices,
because distinct clone classes have precisely the adjacencies of their base
vertices.  Thus the two flag probabilities differ only through the bad event and
rounding error, which proves the claim.
\end{proof}

\begin{lemma}[Positive probability of the planted root]\label{lem:planted-mass}
With the notation above, for all sufficiently large $N$, the finite random
extension of the unlabelled graph $B_N(G,\theta,\lambda)$ to type $\sigma$ gives
probability at least
\[
  c_{k,\lambda}:=
  \begin{cases}
    1, & k=0,\\[1mm]
    \left(\dfrac{\lambda}{2k}\right)^k, & k>0,
  \end{cases}
\]
to the event that the random labelled embedding
$\widehat\theta:\sigma\to B_N(G,\theta,\lambda)$ satisfies
$\widehat\theta(i)\in C_{\theta(i)}$ for all $i=1,\ldots,k$.
\end{lemma}

\begin{proof}
For $k=0$ there is nothing to prove.  Assume $k>0$.  Every ordered choice of
vertices $\widehat\theta(i)\in C_{\theta(i)}$ is an induced embedding of $\sigma$,
because $\theta$ is an induced embedding into $G$ and the blow-up preserves the
adjacencies between distinct base vertices.  For $N$ sufficiently large,
$|C_{\theta(i)}|\ge(\lambda/(2k))N$ for all $i$.  Thus the number of planted
ordered embeddings is at least $(\lambda/(2k))^kN^k$, while the total number of
ordered embeddings of $\sigma$ into $B_N$ is at most $N^k$.
\end{proof}

\subsection{Root-plantability theorem}

\begin{theorem}[Clone-closed classes are root-plantable]\label{thm:clone-root-plantable}
Let $\K$ be a clone-closed hereditary graph class, let
$\T_1=\T_\K$, and let $\sigma$ be a non-degenerate type of $\T_\K$.  Then
\[
  S_\sigma=Q_\sigma.
\]
Consequently, quotient semantics and ensemble semantics are equivalent for every
$f\in\A^\sigma[\T_{\mathrm{Graph}}]$.
\end{theorem}

\begin{proof}
The inclusion $S_\sigma\subseteq Q_\sigma$ is \cref{lem:support-passes-general}.
We prove the reverse inclusion.  Let $\psi\in Q_\sigma$.  We must show that every
neighbourhood of $\psi$ in $X_\sigma$ intersects the support of some
$\Ext_\sigma(\phi_0)$ with $\phi_0\in Q_0$ and $\phi_0(\brak\sigma)>0$.

Let a basic neighbourhood of $\psi$ be specified by flags
$F_1,\ldots,F_s$ and a tolerance $\eps>0$:
\[
  U=
  \{\chi\in X_\sigma:
    |\chi(F_i)-\psi(F_i)|<\eps\text{ for all }i\}.
\]
By the representation theorem in the constrained class $\T_\K$, choose a
convergent sequence of $\sigma$-flags
$(G_t,\theta_t)$ with $G_t\in\K$, $|G_t|\to\infty$, and
$p(\cdot,(G_t,\theta_t))\to\psi$.  Let $m=\max_i |F_i|$.
Choose $0<\lambda\le 1/2$ so small that $C_m\lambda<\eps/10$, where $C_m$ is the
constant from \cref{lem:planted-estimate}.  Then fix one $t$ large enough that
\[
  |p(F_i,(G_t,\theta_t))-\psi(F_i)|<\eps/10
  \quad\text{for all }i,
  \qquad
  \frac{C_m}{|G_t|-|\sigma|}<\eps/10.
\]
Finally take a sequence $N_j\to\infty$ with rounding errors tending to zero and, for all large $j$, satisfying $C_m\operatorname{err}_{N_j}<\eps/10$; put
\[
  B_j:=B_{N_j}(G_t,\theta_t,\lambda).
\]
Since $G_t\in\K$ and $\K$ is clone-closed, every $B_j$ lies in $\K$.  Passing to
a subsequence if necessary, the unlabelled graphs $B_j$ converge to some
$\phi_0\in Q_0$.  The counting argument in \cref{lem:planted-mass} gives a positive lower bound, independent of $j$, on the density of embeddings of $\sigma$ in $B_j$; hence $\phi_0(\brak\sigma)>0$.

Let $Y_\sigma:=[0,1]^{\F^\sigma[\T_{\mathrm{Graph}}]}$, so that
$X_\sigma\subseteq Y_\sigma$ is a closed subspace.  Let $P_j$ be the finite
random-labelling distribution on $Y_\sigma$ obtained by choosing a random embedding
of $\sigma$ into $B_j$ and recording all labelled flag densities.  These measures
converge weakly to $\Ext_\sigma(\phi_0)$, viewed as a probability measure on
$Y_\sigma$ concentrated on $X_\sigma$.

Consider the closed cylinder in the ambient product space
\[
  \widetilde C=
  \{x\in Y_\sigma:
    |x(F_i)-\psi(F_i)|\le\eps/2\text{ for all }i\}.
\]
Put $C:=\widetilde C\cap X_\sigma$.  Then $C\subseteq U$.  For all sufficiently
large $j$, \cref{lem:planted-estimate} implies that every planted root lies
in $\widetilde C$.  By \cref{lem:planted-mass},
\[
  P_j(\widetilde C)\ge c_{|\sigma|,\lambda}>0
\]
for all sufficiently large $j$.  Since $\widetilde C$ is closed in $Y_\sigma$, the
closed-set form of the Portmanteau theorem gives
\[
  \Ext_\sigma(\phi_0)(\widetilde C)
  \ge
  \limsup_{j\to\infty} P_j(\widetilde C)
  \ge c_{|\sigma|,\lambda}>0.
\]
Since $\Ext_\sigma(\phi_0)$ is concentrated on $X_\sigma$, this says
$\Ext_\sigma(\phi_0)(C)>0$.  Hence the support of $\Ext_\sigma(\phi_0)$ intersects
$C\subseteq U$.  Since $U$ was arbitrary, $\psi\in S_\sigma$.
\end{proof}

\begin{corollary}[Clique-free and triangle-free classes]\label{cor:clique-free}
For every $r\ge3$, every non-degenerate $K_r$-free type $\sigma$, and every
$f\in\A^\sigma[\T_{\mathrm{Graph}}]$, quotient semantics and ensemble semantics
are equivalent for the $K_r$-free class.  In particular, the equivalence holds
for the triangle-free class.
\end{corollary}

\section{Degenerate classes are not root-plantable}

The clone-closed hypothesis is not merely technical: the correspondence can fail for
hereditary classes, and it does so for an entire family of them, always for the same
structural reason.  The obstruction is \emph{degeneracy}.  If a class is so sparse
that its dense limits carry no edges, then a typical vertex sees an empty
neighbourhood; yet the class may still contain a graph with an \emph{exceptional}
high-degree vertex, such as the centre of a star.  The labelled edge flag separates the
two.

Throughout this section $\sigma=1$ is the one-vertex type,
$e\in\A^1[\T_{\mathrm{Graph}}]$ is the labelled edge flag (the two-vertex flag in which
the labelled vertex is adjacent to the one unlabelled vertex), and $\rho\in\A^0$ is
the unlabelled edge flag; recall $\avg{e}{1}=\rho$ and $\brak1=1$.

\begin{definition}[Edge-degenerate class]\label{def:edge-degenerate}
A hereditary graph class $\K$ is \emph{edge-degenerate} if every constrained
unlabelled limit has edge density zero, that is, $\phi_0(\rho)=0$ for all
$\phi_0\in Q_0$.  Equivalently, the maximum number of edges of an $n$-vertex graph
in $\K$ is $o(n^2)$.
\end{definition}

\begin{theorem}[Degeneracy obstruction]\label{thm:degenerate-obstruction}
Let $\K$ be a hereditary graph class that is edge-degenerate and contains
arbitrarily large stars.  Then $(\T_{\mathrm{Graph}},\T_\K,\sigma{=}1)$ is not
root-plantable.  Concretely, $f:=-e\in\A^1[\T_{\mathrm{Graph}}]$ satisfies ensemble
semantics but not quotient semantics, and in fact
\[
  S_1\ \subseteq\ \{\chi\in Q_1:\chi(e)=0\}\ \subsetneq\ Q_1 .
\]
\end{theorem}

\begin{proof}
\emph{Ensemble semantics holds for $-e$.}  Let $\phi_0\in Q_0$.  Since $\brak1=1$ we
have $\phi_0(\brak1)=1>0$, so $\Ext_1(\phi_0)$ is defined.  As $\avg{e}{1}=\rho$,
\eqref{eq:extension-expectation} gives
\[
  \E_{\phi^1\sim\Ext_1(\phi_0)}[\phi^1(e)]
  =\frac{\phi_0(\rho)}{\phi_0(\brak1)}=\phi_0(\rho)=0,
\]
using edge-degeneracy.  Because $\phi^1(e)\ge0$ pointwise, $\phi^1(e)=0$ almost
surely, hence $\phi^1(-e)=0\ge0$ almost surely; ensemble semantics holds.  Moreover,
by \cref{lem:support-as} every support point $\chi$ of every $\Ext_1(\phi_0)$
satisfies $\chi(e)=0$, and since $\chi\mapsto\chi(e)$ is continuous this gives the
inclusion $S_1\subseteq\{\chi\in Q_1:\chi(e)=0\}$.

\emph{Quotient semantics fails for $-e$.}  Let $S_N=K_{1,N-1}$ be the star on $N$
vertices, labelled at its centre; by hypothesis $S_N\in\K$ for arbitrarily large $N$.
Every unlabelled vertex is adjacent to the centre, so $p(e,S_N)=1$ for all $N$.  By
compactness a subsequence converges to some $\psi\in Q_1$ with $\psi(e)=1$, whence
$\psi(-e)=-1<0$.  Thus quotient semantics fails, and $\psi\in Q_1$ witnesses
$\{\chi\in Q_1:\chi(e)=0\}\subsetneq Q_1$; in particular $\psi\in Q_1\setminus S_1$,
so $S_1\ne Q_1$.
\end{proof}

\begin{remark}
The hypothesis ``contains arbitrarily large stars'' is only used to exhibit, in the
quotient, a labelled limit of positive labelled edge density; it may be replaced by the
existence of any sequence of labelled graphs in $\K$ whose labelled edge density stays
bounded away from $0$.
\end{remark}

\begin{remark}[The obstruction]
The centre of a star is an exceptional labelled vertex.  Quotient semantics can see
it, because labelled limits may sit at exceptional vertices.  Ensemble semantics
instead labels a constrained unlabelled limit at a \emph{typical} vertex; in a
degenerate class the edge density is zero, so a typical vertex has labelled edge
density zero.  The high-degree vertex cannot be planted with positive probability
without leaving the class: duplicating the centre and one neighbour already produces
a $K_{2,2}$, and more generally a denser local configuration than a degenerate class can
sustain.
\end{remark}

\subsection{Examples}

The hypotheses of \cref{thm:degenerate-obstruction} are met by many familiar
sparse classes.  We first record the original $C_4$-free case, where edge-degeneracy
is the classical K\H{o}v\'ari--S\'os--Tur\'an bound.

Let
\[
  \K_{C_4}:=\{G:\text{$G$ contains no copy of $C_4$ as a not-necessarily-induced subgraph}\}.
\]
This is hereditary; in the induced-subgraph language of flag algebras it is
axiomatised by forbidding the three four-vertex graphs that contain a $C_4$ as a
subgraph: $C_4$, $K_4$ minus an edge, and $K_4$.  It is not clone-closed: the
independent blow-up of $K_2$ into two clone classes of size $2$ is $K_{2,2}=C_4$.

\begin{lemma}[Dense \texorpdfstring{$C_4$}{C4}-free limits have edge density zero]\label{lem:c4-edge-zero}
$\K_{C_4}$ is edge-degenerate: if $\phi_0\in Q_0$ for the $C_4$-free class, then
$\phi_0(\rho)=0$.
\end{lemma}

\begin{proof}
By the representation theorem in the constrained class, $\phi_0$ is the limit
of a convergent sequence of finite $C_4$-free graphs $G_n$ with $|G_n|\to\infty$.
We show that their edge densities tend to zero.

Let $G$ be a $C_4$-free graph on $N$ vertices, with degrees $d(v)$.  Any two
vertices have at most one common neighbour; otherwise two common neighbours and
the two original vertices form a $C_4$.  Therefore
\[
  \sum_{v\in V(G)} \binom{d(v)}{2}
  \le
  \binom{N}{2}.
\]
By convexity,
\[
  N\binom{\overline d}{2}
  \le
  \binom{N}{2},
  \qquad
  \overline d:=\frac1N\sum_v d(v).
\]
Thus $\overline d(\overline d-1)\le N-1$, so
$\overline d\le 1+\sqrt N$.  The edge density is
$\overline d/(N-1)$, which tends to zero.  Hence
$\phi_0(\rho)=0$.
\end{proof}

\begin{corollary}[\texorpdfstring{$C_4$}{C4}-free counterexample]\label{cor:c4-counterexample}
$\K_{C_4}$ is edge-degenerate (\cref{lem:c4-edge-zero}) and contains all stars,
so by \cref{thm:degenerate-obstruction} it is not root-plantable at
$\sigma=1$: the element $-e$ satisfies ensemble semantics but not quotient
semantics.
\end{corollary}

\begin{corollary}[A family of degenerate obstructions]\label{cor:degenerate-family}
Each of the following hereditary classes is edge-degenerate and contains all stars,
hence is not root-plantable at the one-vertex type:
\begin{enumerate}[label=\textup{(\roman*)}, leftmargin=2.5em]
  \item $K_{s,t}$-subgraph-free graphs for $2\le s\le t$, edge-degenerate by the
        K\H{o}v\'ari--S\'os--Tur\'an bound $\mathrm{ex}(n,K_{s,t})=O(n^{2-1/s})$
        \cite{KovariSosTuran} (a star contains no $K_{2,2}$);
  \item $C_{2k}$-subgraph-free graphs for $k\ge2$, edge-degenerate by the
        Bondy--Simonovits bound $\mathrm{ex}(n,C_{2k})=O(n^{1+1/k})$
        \cite{BondySimonovits};
  \item forests, and more generally any class of bounded average degree containing
        all stars ($O(n)$ edges);
  \item planar graphs ($\le 3n-6$ edges).
\end{enumerate}
\end{corollary}

\subsection{Complementation: a dense obstruction}

One might hope that the failures above are an artefact of \emph{sparsity}, and that
density rescues root-plantability.  It does not.  Root-plantability is preserved by
complementation, which immediately turns each sparse obstruction into a dense one.

\begin{lemma}[Complementation invariance]\label{lem:complementation}
For a hereditary graph class $\K$ let $\overline\K=\{\overline G:G\in\K\}$, which is
again hereditary, and for a type $\sigma$ let $\overline\sigma$ be its complement.
Then $(\T_{\mathrm{Graph}},\T_\K,\sigma)$ is root-plantable if and only if
$(\T_{\mathrm{Graph}},\T_{\overline\K},\overline\sigma)$ is.  For the one-vertex type
$\sigma=1=\overline\sigma$, the class $\K$ is root-plantable at $1$ if and only if
$\overline\K$ is.
\end{lemma}

\begin{proof}
Vertex-wise complementation $G\mapsto\overline G$ is an involution on finite graphs
that takes induced-subgraph densities to those of the complements, hence extends to
an isomorphism of flag algebras
$\A^\sigma[\T_{\mathrm{Graph}}]\cong\A^{\overline\sigma}[\T_{\mathrm{Graph}}]$ carrying
$\K$-flags to $\overline\K$-flags.  It therefore induces a homeomorphism
$X_\sigma\cong X_{\overline\sigma}$ taking $Q_\sigma$ for $\K$ onto $Q_{\overline\sigma}$
for $\overline\K$.  Complementation also commutes with the random labelling: a uniformly
random embedding of $\sigma$ into $G$ is the same embedding of $\overline\sigma$ into
$\overline G$, so it carries each $\Ext_\sigma(\phi_0)$ to the corresponding
extension for $\overline\K$, and hence $S_\sigma$ for $\K$ onto $S_{\overline\sigma}$
for $\overline\K$.  Thus $S_\sigma=Q_\sigma$ holds for $\K$ if and only if it holds
for $\overline\K$.
\end{proof}

Call $\K$ \emph{co-edge-degenerate} if $\overline\K$ is edge-degenerate; equivalently
every $\phi_0\in Q_0$ has $\phi_0(\rho)=1$, equivalently every $n$-vertex graph in
$\K$ has at least $\binom n2-o(n^2)$ edges.  Such classes are as dense as possible,
with edge density tending to $1$, yet they fail just as the sparse ones do.

\begin{corollary}[A dense obstruction]\label{cor:codegenerate}
Let $\K$ be hereditary, co-edge-degenerate, and contain all co-stars
$\overline{K_{1,N-1}}=K_{N-1}\sqcup K_1$.  Then $\K$ is not root-plantable at
$\sigma=1$, witnessed by $f:=e-\one_1$ (with $\one_1$ the unit of $\A^1$).  In particular the complements of the classes in
\cref{cor:degenerate-family}---complement-of-$C_4$-free,
complement-of-$K_{s,t}$-free, complement-of-$C_{2k}$-free, and
complement-of-planar graphs---are dense hereditary classes that are not
root-plantable.
\end{corollary}

\begin{proof}
By \cref{lem:complementation} it suffices that $\overline\K$ be not
root-plantable at $1$, and this is \cref{thm:degenerate-obstruction}:
$\overline\K$ is edge-degenerate and contains all stars $K_{1,N-1}$.  The witness
$-e$ for $\overline\K$ pulls back to $-(\one_1-e)=e-\one_1$ for $\K$.  Directly:
$\avg{e}{1}=\rho$ gives $\E[\phi^1(e)]=\phi_0(\rho)=1$, and $\phi^1(e)\le1$ forces
$\phi^1(e)=1$ almost surely, so $e-\one_1$ satisfies ensemble semantics; while the
co-star labelled at its isolated vertex gives a $\psi\in Q_1$ with $\psi(e)=0$, whence
$\psi(e-\one_1)=-1<0$ and quotient semantics fails.
\end{proof}

\begin{remark}[The dichotomy: density is not the dividing line]
\cref{thm:degenerate-obstruction} and \cref{cor:codegenerate}
together show that root-plantability is governed neither by heredity nor by density:
edge-degenerate classes (labelled edge density pinned to $0$) and co-edge-degenerate
classes (pinned to $1$) both fail, for the same reason.  What separates a typical
vertex from the exceptional one is that the labelled edge density is \emph{pinned} to a
boundary value of $[0,1]$ across all constrained unlabelled limits, while the
constrained class still contains a labelled limit attaining the opposite boundary.  The
governing invariant is this boundary pinning, not sparsity or density.

This refines, rather than closes, the problem of characterising root-plantability.
The $C_5$-free graphs show why the question must be asked at a fixed type.  At the
one-vertex type their edge density is \emph{not} pinned---the class contains graphs of
every density in $[0,\tfrac12]$, so no boundary-pinning obstruction applies---and a
local root-cloning repair proves root-plantability there.  The same repair also handles
the two-root non-edge type by cloning the two non-adjacent roots simultaneously.  At the
two-root edge type, however, the edge-rooted triangle flag is pinned to $0$ on random
edge extensions, while a book graph gives a quotient edge with common-neighbour density
$1$.  Thus the full all-types $C_5$-free conjecture is false, even though the positive
one-root and non-edge cases remain useful models for planting beyond clone-closure.
\end{remark}

\subsection{The general obstruction and a conjecture}

Both obstructions are instances of one mechanism: a labelled density pinned almost
surely to a single value across all constrained limits, while the quotient escapes
it.

\begin{theorem}[Pinning obstruction]\label{thm:pinning}
Let $\sigma$ be a non-degenerate type, let $g$ be a $\sigma$-flag, and let
$c\in[0,1]$.  Suppose the labelled density of $g$ is almost surely pinned to $c$ across
all constrained unlabelled limits: for every $\phi_0\in Q_0$ with
$\phi_0(\brak\sigma)>0$, one has $\phi^\sigma(g)=c$ almost surely under
$\Ext_\sigma(\phi_0)$.  If the constrained class contains a labelled limit attaining a different
value---some $\psi\in Q_\sigma$ with $\psi(g)\ne c$---then $(\T_0,\T_1,\sigma)$ is not
root-plantable.
\end{theorem}

\begin{proof}
Both $g-c\,\one_\sigma$ and $c\,\one_\sigma-g$ have $\phi^\sigma$-value $0$ almost surely
under every admissible $\Ext_\sigma(\phi_0)$, so each is ensemble-nonnegative and,
by \cref{lem:support-as} and continuity, nonnegative on $S_\sigma$.  Since
$\psi(g)\ne c$, one of the two is strictly negative at $\psi$; that
$\psi\in Q_\sigma$ therefore lies outside $S_\sigma$, so $S_\sigma\ne Q_\sigma$.
\end{proof}

\begin{remark}[Endpoints are automatic and checkable]
At the endpoints $c\in\{0,1\}$ the almost-sure pinning follows from a mere
expectation condition, since a $[0,1]$-valued random variable with mean $0$
(respectively $1$) is almost surely $0$ (respectively $1$).  Concretely, by
\eqref{eq:extension-expectation}, $\phi_0(\avg{g}{\sigma})=0$ for all admissible
$\phi_0$ gives the case $c=0$ (witness $-g$), and
$\phi_0(\avg{g}{\sigma})=\phi_0(\brak\sigma)$ gives $c=1$ (witness $g-\one_\sigma$).  The
degeneracy obstruction (\cref{thm:degenerate-obstruction}) and its dense
counterpart (\cref{cor:codegenerate}) are exactly these, with $\sigma=1$ and
$g=e$.  An interior value $c\in(0,1)$ would instead require the almost-sure pinning as a
genuine hypothesis---but for a subgraph-forbidden class no interior value is ever pinned,
as \cref{subsec:boundary} shows.
\end{remark}

\begin{conjecture}[Characterisation; tentative]\label{conj:characterisation}
$(\T_{\mathrm{Graph}},\T_\K,\sigma)$ is root-plantable if and only if it admits no
pinning obstruction: there is no $\sigma$-flag $g$ and constant $c$ with
$\phi^\sigma(g)=c$ almost surely for all constrained unlabelled limits while some
$\psi\in Q_\sigma$ has $\psi(g)\ne c$.
\end{conjecture}

The forward implication is \cref{thm:pinning}; the conjecture asserts that the
absence of such single-flag pinning is already sufficient.  It is tentative:
root-plantability could conceivably fail through a \emph{joint} constraint on several
flags with no single flag pinned, in which case a richer notion of obstruction would
be needed.

\subsection{Pinning is always to a boundary value}\label{subsec:boundary}

The obstructions exhibited above are all at the boundary, $c\in\{0,1\}$, and the remark
after \cref{thm:pinning} observed that an interior value would require the almost-sure
constancy as a genuine hypothesis.  We now show that for a subgraph-forbidden class no
interior value is ever pinned: every pinning obstruction is a boundary one.  This matters
because, as the next section shows, boundary obstructions are inert for density bounds.

\begin{theorem}[No interior pinning]\label{thm:no-interior}
Let $\K$ be defined by forbidding a family of graphs as (not necessarily induced)
subgraphs, and let $\sigma$ be any non-degenerate type.  If a $\sigma$-flag $g$ is pinned
to $c$ on $S_\sigma$, then $c\in\{0,1\}$.
\end{theorem}

\begin{proof}
It is enough to produce a single point $o_\sigma\in S_\sigma$ that is $\{0,1\}$-valued on
every flag: pinning forces $g\equiv c$ on $S_\sigma$, so then $c=o_\sigma(g)\in\{0,1\}$.

As $\sigma$ is non-degenerate, some $\phi_0\in Q_0$ has positive $\sigma$-density; let $W$
be a graphon representing it.  For $\lambda\in(0,1]$ put $W_\lambda:=\lambda W$.  Then
$W_\lambda\le W$ pointwise, and the density of any subgraph is monotone in the edge
weights, so every forbidden subgraph still has density $0$ in $W_\lambda$; hence
$W_\lambda$ is $F$-free and its limit $\phi_0^{\lambda}$ lies in $Q_0$.  Its
$\sigma$-density is positive for each $\lambda>0$, because the integrand defining the
induced $\sigma$-density is positive on the positive-measure set of tuples on which $W$
realises $\sigma$; so $\phi_0^{\lambda}$ can be labelled by $\sigma$.

Every edge weight of $W_\lambda$ is at most $\lambda$, so under
$\Ext_\sigma(\phi_0^{\lambda})$ the probability that a given further vertex attaches to
a labelled vertex, or that two further vertices are adjacent, is at most $\lambda$.  Hence as
$\lambda\to0$ the labelled density of a $\sigma$-flag $g$ tends to $1$ when $g$ extends
$\sigma$ by pairwise non-adjacent vertices that are also non-adjacent to $\sigma$, and to
$0$ otherwise.  These limiting values are $\{0,1\}$, so $\Ext_\sigma(\phi_0^{\lambda})$
concentrates at the labelled limit $o_\sigma$---the type $\sigma$ together with an isolated
edgeless cloud---which therefore lies in $S_\sigma$ and is $\{0,1\}$-valued on every flag.
\end{proof}

\begin{remark}[Scope]
\Cref{thm:no-interior} is uniform in $\sigma$, with no restriction on the type, and
subsumes the one-vertex, edge, and clique cases at once.  The scaled graphons
$\lambda W$ supply, as $\lambda\to0$, a deterministic labelled view of $\sigma$ from which
the boundary value is read off.  No single $\K$-limit of positive $\sigma$-density has
isolated $\sigma$-copies; but letting the $\sigma$-density tend to zero \emph{along} the
family while staying positive realises such a view in the closure $S_\sigma$, exactly as
the degenerate corner of a complete bipartite family does for the edge type.  By
complementation (\cref{lem:complementation}) the same holds for every class whose
complement is subgraph-forbidden, including the dense co-degenerate ones.  The
consequences for density bounds are drawn in \cref{sec:empty-type}.
\end{remark}

\section{The gap is invisible to density bounds}\label{sec:empty-type}

The obstructions above all occur at the \emph{labelled} type $\sigma=1$: the witnesses
$-e$ and $e-\one_1$, and the pinning witness $g-c\,\one_\sigma$ of
\cref{thm:pinning}, are elements of $\A^\sigma[\T_0]$ with $\sigma\ne\emptyset$.  A
flag-algebra \emph{application}, by contrast, outputs a density inequality in the
empty-type algebra $\A^0[\T_1]$---a statement $f\ge0$ on $Q_0$, such as a Mantel or
Tur\'an bound.  One should therefore ask whether the incompleteness just exhibited can
ever affect such a bound.  It cannot, for two independent reasons: at the empty type the
two semantics coincide, and the witnesses that expose the gap unlabel to zero.

\begin{proposition}[Empty-type collapse]\label{prop:empty-type}
For the empty type, $\Ext_\emptyset(\phi_0)=\delta_{\phi_0}$ for every $\phi_0\in Q_0$;
hence $S_\emptyset=Q_\emptyset=Q_0$.  Thus $(\T_0,\T_1,\emptyset)$ is always
root-plantable, and quotient and ensemble semantics coincide for every $f\in\A^0[\T_0]$.
\end{proposition}

\begin{proof}
The empty type admits a unique embedding into any graph, so the random labelling carries
no information; we read this off \eqref{eq:extension-expectation}.  For
$\sigma=\emptyset$ the unlabelling operator is the identity and $\brak\emptyset=\one$, so
$\phi_0(\brak\emptyset)=1$ and $\mu:=\Ext_\emptyset(\phi_0)$ satisfies
$\E_{\psi\sim\mu}[\psi(g)]=\phi_0(g)$ for all $g\in\A^0$.  Applying this to $g$ and to
$g^2$ and using that $\phi_0$ is multiplicative,
\[
  \E_{\psi\sim\mu}\bigl[(\psi(g)-\phi_0(g))^2\bigr]=\phi_0(g^2)-\phi_0(g)^2=0,
\]
so $\psi(g)=\phi_0(g)$ for $\mu$-almost every $\psi$, for each of the countably many
$g\in\A^0$.  Since flag densities separate the points of $X_0$, this forces
$\mu=\delta_{\phi_0}$, whence
$S_\emptyset=\overline{\bigcup_{\phi_0}\{\phi_0\}}=\overline{Q_0}=Q_0$.  The coincidence
of the two semantics now follows from \cref{thm:support-criterion}.
\end{proof}

\begin{corollary}\label{cor:confined}
No density bound in $\A^0[\T_1]$ is ensemble-true but quotient-false; every failure of
the quotient/ensemble correspondence occurs at a labelled type $\sigma\ne\emptyset$.  In
particular the witnesses $-e$ and $e-\one_1$ are inequalities at $\sigma=1$, not density
bounds.
\end{corollary}

\Cref{cor:confined} settles the question of \emph{truth}.  A flag-algebra proof, though,
reaches its empty-type conclusion through labelled multipliers, so one must still ask
whether the labelled gap can affect a \emph{certificate}.  The method writes
\[
  f-\sum_i\avg{s_i}{\sigma_i}=(\text{a manifestly non-negative combination of flags}),
\]
each $s_i$ a sum of squares in $\A^{\sigma_i}[\T_1]$.  Because $\avg{s_i}{\sigma_i}$
integrates $s_i$ against $\Ext_{\sigma_i}(\phi_0)$, a measure supported on
$S_{\sigma_i}$, non-negativity of $s_i$ on $S_{\sigma_i}$ already suffices---weaker than
the sum-of-squares condition, which is non-negativity on all of $Q_{\sigma_i}$.  An
\emph{ensemble-relaxed} method could therefore use any multiplier non-negative on
$S_\sigma$.  Write, as cones inside $\A^0[\T_1]$,
\[
  \C^{\mathrm{quot}}_\sigma=\{\avg{s}{\sigma}:s\text{ a sum of squares}\},
  \qquad
  \C^{\mathrm{ens}}_\sigma=\{\avg{s}{\sigma}:s\ge0\text{ on }S_\sigma\},
\]
for the empty-type contributions the type $\sigma$ can make to the two methods, so that
$\C^{\mathrm{quot}}_\sigma\subseteq\C^{\mathrm{ens}}_\sigma$ always.  The next two results
show the witnesses gain nothing here, and that for the obstructions exhibited above the
two cones in fact coincide.

\begin{proposition}[The vanishing ideal unlabels to zero]\label{prop:ideal-zero}
Let $\mathcal I_\sigma=\{a\in\A^\sigma[\T_1]:\psi(a)=0\text{ for all }\psi\in S_\sigma\}$,
where each $\psi\in S_\sigma\subseteq Q_\sigma$ is evaluated on $\A^\sigma[\T_1]$ through
the unique homomorphism $\bar\psi\in\Homp(\A^\sigma[\T_1],\R)$ it factors as,
$\psi=\bar\psi\circ\qmap_\sigma$.
Then $\avg{a}{\sigma}=0$ in $\A^0[\T_1]$ for every $a\in\mathcal I_\sigma$.  In
particular every pinning witness $g-c\,\one_\sigma$, and every flag-multiple
$(g-c\,\one_\sigma)\,h$, unlabels to zero.  Consequently $\avg{\cdot}{\sigma}$ depends
only on the restriction $s|_{S_\sigma}$, so
$\C^{\mathrm{ens}}_\sigma\supsetneq\C^{\mathrm{quot}}_\sigma$ requires some $s\ge0$ on
$S_\sigma$ whose restriction $s|_{S_\sigma}$ is the restriction of no sum of squares.
\end{proposition}

\begin{proof}
For $\phi_0\in Q_0$ with $\phi_0(\brak\sigma)>0$,
$\phi_0(\avg{a}{\sigma})=\phi_0(\brak\sigma)\,\E_{\psi\sim\Ext_\sigma(\phi_0)}[\psi(a)]$
by \eqref{eq:extension-expectation}, and the expectation vanishes because the measure is
supported on $S_\sigma$ and $a=0$ there.  If $\phi_0(\brak\sigma)=0$, then
$\avg{a}{\sigma}$ is a combination of unlabelled graphs each containing a copy of
$\sigma$, so each has $\phi_0$-density $0$ by monotonicity and again
$\phi_0(\avg{a}{\sigma})=0$.  Thus $\avg{a}{\sigma}$ vanishes on $Q_0$ and so is the zero
element.  A pinning witness lies in $\mathcal I_\sigma$ by definition, and
$\mathcal I_\sigma$ is an ideal; the final clause holds because
$s-s'\in\mathcal I_\sigma$ whenever $s,s'$ agree on $S_\sigma$.
\end{proof}

\begin{proposition}[Degenerate roots collapse to a point]\label{prop:single-point}
If $\K$ is edge-degenerate, or co-edge-degenerate, then $S_1$ is a single point.
Consequently $\C^{\mathrm{ens}}_1=\C^{\mathrm{quot}}_1=\R_{\ge0}\,\one_0$: at the
one-vertex type the ensemble relaxation certifies exactly the density bounds the
quotient method does.  The degeneracy and co-degeneracy counterexamples
(\cref{cor:c4-counterexample,cor:degenerate-family,cor:codegenerate}) therefore cost
nothing for density bounds.
\end{proposition}

\begin{proof}
Take $\K$ edge-degenerate; the co-degenerate case follows by complementation
(\cref{lem:complementation}).  Let $F$ be any $1$-flag whose underlying graph $M$ has an
edge.  A random $|M|$-subset of an $n$-vertex graph spans, in expectation,
$\binom{|M|}{2}$ times the edge density many edges, so the induced density obeys
$p(M,G)\le\binom{|M|}{2}\rho(G)$; in the limit
$\phi_0(M)\le\binom{|M|}2\phi_0(\rho)=0$ for every $\phi_0\in Q_0$.  As $\avg{F}{1}$ is a
non-negative multiple of $M$, we get
$\E_{\phi^1\sim\Ext_1(\phi_0)}[\phi^1(F)]=\phi_0(\avg{F}{1})=0$, and since
$\phi^1(F)\ge0$, $\phi^1(F)=0$ almost surely.  Thus every support point of every
$\Ext_1(\phi_0)$ assigns density $0$ to every edge-containing $1$-flag; exactly one
homomorphism $o\in X_1$ does so (the labelled empty-graph limit, the others being
determined by the chain relations), so $S_1=\{o\}$.  On a single point every element
non-negative there agrees with the constant $s(o)\,\one_1$, a sum of squares, so
$\avg{s}{1}=s(o)\,\brak1=s(o)\,\one_0$ for every such $s$; both cones equal
$\R_{\ge0}\,\one_0$.
\end{proof}

\begin{remark}[What a practically relevant gap would require]
\Cref{prop:single-point} shows the gap is not merely witness-inert but completely inert
for the exhibited classes: a degenerate labelling carries no information, so its only
contribution to a density bound is an additive constant, in either semantics.  A gap
that \emph{does} reach a density bound therefore requires, by \cref{prop:ideal-zero}, a
\emph{Positivstellensatz gap on the support variety}: a non-root-plantable type whose
$S_\sigma$ is rich enough to carry a non-negative element that is not a restricted sum of
squares.  Boundary pinning to a single point supplies none, and by
\cref{thm:no-interior} pinning is \emph{always} to a boundary value---so every pinning
obstruction, exhibited here or not, is boundary.  \Cref{conj:characterisation} predicts
that the only obstructions are pinning; if so, and if the boundary support sets remain
sum-of-squares-saturated, the quotient/ensemble gap---though genuine at labelled
types---would never affect a density bound.  The $C_5$-free case of the next section
illustrates both possibilities at once: the one-root and non-edge-rooted supports are
rich but plantable, while the edge-rooted support has a boundary pinning obstruction
coming from common neighbours of an exceptional book edge.
\end{remark}

\section{The \texorpdfstring{$C_5$}{C5}-free split}

The $C_5$-free graphs are dense and not clone-closed, so they are the first natural
place to test root-plantability beyond the blow-up theorem.  They split sharply by
type.  At the one-vertex type, the elementary pinning obstructions disappear and a local
planting construction works.  At the two-root non-edge type, a simultaneous version of
the same construction works.  At the two-root edge type, a book-edge construction gives a
pinning obstruction.  Throughout, $C_5$-free again means containing no $C_5$ as a
not-necessarily-induced subgraph (the densest such graphs being complete bipartite).

The key fact is a strong local constraint.

\begin{lemma}[Neighbourhoods in $C_5$-free graphs]\label{lem:c5-nbhd}
If $G$ is $C_5$-free and $v\in V(G)$, then $G[N(v)]$ contains no path on four
vertices as a subgraph.  Consequently every connected component of $G[N(v)]$ is
a star or a triangle, and in particular
\[
  e(G[N(v)])\le |N(v)|.
\]
\end{lemma}

\begin{proof}
A path $a$-$x$-$y$-$b$ inside $N(v)$ together with $v$ forms the $5$-cycle
$v$-$a$-$x$-$y$-$b$-$v$.  Thus $G[N(v)]$ is $P_4$-subgraph-free.

Now let $C$ be a connected component of a graph with no $P_4$ subgraph.  If $C$
is acyclic, then its diameter is at most $2$, hence it is a star.  If $C$
contains a cycle, no cycle has length at least $4$, and a triangle cannot have
any further vertex in the same component: a shortest path from such a vertex to
the triangle would contain a $P_4$.  Thus $C$ is a triangle.  Each component has
at most as many edges as vertices, giving the displayed bound.
\end{proof}

\begin{corollary}[No pinning obstruction for $C_5$-free]\label{cor:c5-no-pin}
The $C_5$-free class has no pinning obstruction at $\sigma=1$ from the edge or the
labelled triangle flag.
\end{corollary}

\begin{proof}
A disjoint union of stars and triangles on at most $\deg(v)$ vertices has at most
$\deg(v)$ edges, so the number of triangles through $v$ is at most $\deg(v)\le n$,
and the labelled triangle density at every vertex is $O(1/n)$.  Hence the
labelled triangle flag has density $0$ on \emph{all} of $Q_1$: it is pinned at $0$, but
the quotient attains only $0$, so no obstruction arises.  The edge flag is not
pinned at all, as $C_5$-free graphs realise every edge density in $[0,\tfrac12]$
(the empty graph and the complete bipartite graphs).
\end{proof}

\subsection{The one-root planting theorem}

The one-root case is the first example where root-plantability holds without
clone-closure.  Ordinary false-twin duplication of a vertex may create $C_5$s:
if the root sits in a triangle, cloning the root and one neighbour creates a
five-cycle.  The repair is local.  Clone the root into a small positive-density
independent set, but first delete all old edges inside its neighbourhood.  The
deleted set is only linear in the number of vertices by \cref{lem:c5-nbhd}, so
it is invisible to fixed rooted densities.

\begin{definition}[One-root planting]\label{def:c5-one-root-planting}
Let $(G,r)$ be a rooted graph and let $L\ge1$.  The planted graph
$P_L(G,r)$ is defined as follows.  Remove $r$, add an independent set $R$ of
$L$ new vertices, and keep the old vertex set
\[
  U:=V(G)\setminus\{r\}.
\]
Inside $U$, keep every edge of $G-r$ except the edges with both endpoints in
$N_G(r)$.  Finally, join every vertex of $R$ to every vertex of $N_G(r)$, and
to no other vertex of $U$.
\end{definition}

\begin{lemma}[The planted graph is still $C_5$-free]\label{lem:c5-planting-free}
If $G$ is $C_5$-free, then $P_L(G,r)$ is $C_5$-free for every $L\ge1$.
\end{lemma}

\begin{proof}
Put $H=P_L(G,r)$ and let $R$ be the planted root cluster.  Suppose, for a
contradiction, that $H$ contains a $5$-cycle $C$.

If $C$ uses no vertex of $R$, then all its edges are old edges of $G-r$, so it
is a $C_5$ in $G$.

If $C$ uses exactly one vertex $x\in R$, write the cycle as
\[
  x-a-b-c-d-x.
\]
The edges incident with $x$ force $a,d\in N_G(r)$, and the three middle edges
are old edges of $G$.  Hence
\[
  r-a-b-c-d-r
\]
is a $C_5$ in $G$, again impossible.

Finally suppose $C$ uses at least two vertices of $R$.  Since $R$ is
independent and the independence number of $C_5$ is $2$, it uses exactly two,
say $x,y\in R$.  On one of the two arcs between $x$ and $y$ in the cycle there
are two consecutive old vertices, say $b,c$.  Both are adjacent in $H$ to a
vertex of $R$, so both lie in $N_G(r)$; but the edge $bc$ would then be an edge
inside $N_G(r)$, and all such edges were deleted in the construction of $H$.
This contradiction proves the lemma.
\end{proof}

\begin{lemma}[Rooted density approximation]\label{lem:c5-density-approx}
Fix $m\ge1$.  There is a constant $C_m$ such that the following holds.  Let
$(G,r)$ be a rooted $C_5$-free graph on $n$ vertices, let
$H=P_L(G,r)$ with $1\le L\le \lambda n$, and root $H$ at any vertex
$x\in R$.  For every one-rooted flag $F$ with at most $m$ vertices,
\[
  \left|p(F,(H,x))-p(F,(G,r))\right|
  \le C_m\left(\lambda+\frac1n\right).
\]
\end{lemma}

\begin{proof}
After enlarging $C_m$ we may ignore the finitely many cases with $n$ bounded in
terms of $m$, since the left-hand side is always at most $1$.  Assume therefore
that all denominators below are non-zero and comparable to the displayed powers
of $n$.

Let $\ell=|F|\le m$.  The density $p(F,(H,x))$ is computed by sampling
$\ell-1$ further vertices from $H\setminus\{x\}$.  Couple this with a uniform
sample from $U=V(G)\setminus\{r\}$.

There are two bad events.  First, the sample in $H$ may hit another vertex of
the planted cluster $R$; this has probability at most
\[
  \frac{(\ell-1)(L-1)}{n+L-2}=O_m(\lambda+1/n).
\]
Second, after the sample has landed in $U$, it may contain both endpoints of an
edge deleted from $G[N_G(r)]$.  By \cref{lem:c5-nbhd}, the number of deleted
edges is at most $n$, so this probability is at most
\[
  \binom{\ell-1}{2}\frac{n}{\binom{n-1}{2}}
  =O_m(1/n).
\]
Outside these bad events, the rooted induced graph seen from $x$ in $H$ is
identical to the rooted induced graph seen from $r$ in $G$: the root has the
same neighbours in $U$, and no edge among the sampled old vertices has been
changed.  Enlarging the implicit constant gives the stated bound.
\end{proof}

\begin{theorem}[One-root $C_5$-free root-plantability]\label{thm:c5-one-root}
The $C_5$-free graph class is root-plantable at the one-vertex type:
\[
  S_1=Q_1.
\]
\end{theorem}

\begin{proof}
The inclusion $S_1\subseteq Q_1$ is the general support-passing statement for
hereditary constrained classes.  We prove $Q_1\subseteq S_1$.

Let $\psi\in Q_1$, and let $U$ be a basic neighbourhood of $\psi$, specified by
one-rooted flags $F_1,\ldots,F_s$ of size at most $m$ and a tolerance
$\eps>0$:
\[
  U=
  \{\chi\in Q_1:
    |\chi(F_i)-\psi(F_i)|<\eps\text{ for all }i\}.
\]
Choose a convergent sequence of rooted $C_5$-free graphs
$(G_t,r_t)$ with $|G_t|\to\infty$ and
$p(\,\cdot\,,(G_t,r_t))\to\psi$.

Pick $\lambda>0$ so small that $C_m\lambda<\eps/4$, where $C_m$ is the constant
from \cref{lem:c5-density-approx}.  For each $t$, put
\[
  L_t=\max(1,\lfloor\lambda |G_t|\rfloor),
  \qquad
  H_t:=P_{L_t}(G_t,r_t).
\]
By \cref{lem:c5-planting-free}, every $H_t$ is $C_5$-free.  Passing to a
subsequence, assume the unlabelled graphs $H_t$ converge to some
$\phi_0\in Q_0$.

Let $R_t$ be the planted cluster in $H_t$.  The probability that a uniformly
random root of $H_t$ lies in $R_t$ tends to
\[
  \delta=\frac{\lambda}{1+\lambda}>0.
\]
Moreover, for every $x\in R_t$, \cref{lem:c5-density-approx} and the convergence
of $(G_t,r_t)$ to $\psi$ imply, for all sufficiently large $t$,
\[
  |p(F_i,(H_t,x))-\psi(F_i)|\le\eps/2
  \qquad(1\le i\le s).
\]
Let
\[
  Y_1:=[0,1]^{\F^{1}[\T_{\mathrm{Graph}}]},
  \qquad
  \widetilde C=
  \{x\in Y_1:
    |x(F_i)-\psi(F_i)|\le\eps/2\text{ for all }i\}.
\]
This is a closed cylinder in the ambient product space.  Its intersection with
$X_1$ is contained in $U$.  The finite random-root distribution of $H_t$,
viewed as a measure on $Y_1$ by recording all rooted flag densities, assigns
mass at least $\delta/2$ to $\widetilde C$ for all large $t$.

By Razborov's random-extension theorem, these finite random-root distributions
converge weakly to $\Ext_1(\phi_0)$ as measures on $Y_1$.  The closed-set form
of the Portmanteau theorem gives
\[
  \Ext_1(\phi_0)(\widetilde C)\ge\delta/2>0.
\]
Since $\Ext_1(\phi_0)$ is concentrated on $X_1$, its support intersects
$\widetilde C\cap X_1\subseteq U$.  Since $U$ was an arbitrary neighbourhood of
$\psi$, we have $\psi\in S_1$.
\end{proof}

\begin{remark}
\Cref{thm:c5-one-root} is the promised local replacement for the clone-closed
blow-up theorem in one rooted variable.  It uses two facts special to one root:
the only dangerous new $C_5$s created by cloning the root pass through an old
edge inside $N_G(r)$, and those old neighbourhood edges have density
$O(1/|G|)$ in rooted sampling.  For larger types one would have to control
edges inside and between several root-neighbourhood regions.  The next result
shows that this can still be done for two non-adjacent roots; the following
subsection shows that it cannot work in general at the edge type.
\end{remark}

\subsection{The two-root non-edge type}

Let $\eta$ be the two-vertex non-edge type.  The same idea works for $\eta$:
clone both roots into two independent clusters, leave no edges between the
clusters, and delete the old edges inside each root-neighbourhood.  Mixed
five-cycles using one clone of each root project back to genuine $C_5$s in the
original graph; cycles using two clones from the same root are killed by the
neighbourhood-edge deletion.

\begin{definition}[Two-root non-edge planting]\label{def:c5-nonedge-planting}
Let $(G,r,s)$ be a rooted graph with $rs\notin E(G)$, and let $L\ge1$.  The
planted graph $P_L(G,r,s)$ is defined as follows.  Remove $r,s$, add two
disjoint independent sets $R$ and $S$ of size $L$, and keep the old vertex set
\[
  U:=V(G)\setminus\{r,s\}.
\]
There are no edges inside $R$, inside $S$, or between $R$ and $S$.  Inside
$U$, keep every edge of $G-\{r,s\}$ except the edges with both endpoints in
$N_G(r)$ or both endpoints in $N_G(s)$.  Finally, join every vertex of $R$ to
every vertex of $N_G(r)$, and every vertex of $S$ to every vertex of $N_G(s)$,
with no other new edges.
\end{definition}

\begin{lemma}[The non-edge planted graph is $C_5$-free]\label{lem:c5-nonedge-planting-free}
If $G$ is $C_5$-free and $rs\notin E(G)$, then $P_L(G,r,s)$ is $C_5$-free for
every $L\ge1$.
\end{lemma}

\begin{proof}
Put $H=P_L(G,r,s)$ and let $R,S$ be the planted clusters.  The set
$R\cup S$ is independent, and the independence number of $C_5$ is $2$, so any
$5$-cycle in $H$ uses at most two planted vertices.

If a $5$-cycle uses no planted vertex, all its edges are old edges of
$G-\{r,s\}$, hence it is a $C_5$ in $G$.

If it uses exactly one planted vertex, say $x\in R$, write the cycle as
\[
  x-a-b-c-d-x.
\]
Then $a,d\in N_G(r)$ and the middle edges are old edges of $G$, so
$r-a-b-c-d-r$ is a $C_5$ in $G$.  The case $x\in S$ is identical.

If it uses two planted vertices from the same cluster, say $x,y\in R$, then on
the length-three arc between $x$ and $y$ there are two consecutive old vertices
$b,c$.  Both are adjacent to a vertex of $R$, hence both lie in $N_G(r)$, but
the edge $bc$ would have been deleted.  The same argument applies to two
vertices of $S$.

It remains to consider one planted vertex $x\in R$ and one planted vertex
$y\in S$.  Since $x$ and $y$ are non-adjacent, replacing $x$ by $r$ and $y$ by
$s$ in the cycle preserves all five cycle edges: neighbours of $x$ lie in
$N_G(r)$, neighbours of $y$ lie in $N_G(s)$, and all edges among the three old
vertices are old edges of $G$.  The replacement therefore gives a $C_5$ in
$G$, contradicting the hypothesis.
\end{proof}

\begin{lemma}[Non-edge rooted density approximation]\label{lem:c5-nonedge-density-approx}
Fix $m\ge2$.  There is a constant $C_m$ such that the following holds.  Let
$(G,r,s)$ be a $C_5$-free graph rooted at a non-edge, with $n$ vertices, let
$H=P_L(G,r,s)$ with $1\le L\le\lambda n$, and root $H$ at any ordered pair
$(x,y)\in R\times S$.  For every $\eta$-flag $F$ with at most $m$ vertices,
\[
  \left|p(F,(H,x,y))-p(F,(G,r,s))\right|
  \le C_m\left(\lambda+\frac1n\right).
\]
\end{lemma}

\begin{proof}
As in \cref{lem:c5-density-approx}, bounded $n$ can be absorbed into the
constant $C_m$, so assume all denominator estimates below are valid.  Let
$\ell=|F|\le m$.  The density $p(F,(H,x,y))$ is computed by sampling
$\ell-2$ further vertices from $H\setminus\{x,y\}$.  Couple this with a uniform
sample from $U=V(G)\setminus\{r,s\}$.

The sample in $H$ may hit another planted vertex in $R\cup S$; this has
probability $O_m(\lambda+1/n)$.  Once the sample lies in $U$, the only remaining
difference from $(G,r,s)$ is that edges inside $N_G(r)$ or inside $N_G(s)$ have
been deleted.  By \cref{lem:c5-nbhd}, the number of such deleted edges is at
most $2n$, so the probability that the sample contains the endpoints of one of
them is $O_m(1/n)$.  Outside these bad events, the rooted induced graph seen
from $(x,y)$ in $H$ is identical to the rooted induced graph seen from
$(r,s)$ in $G$: the roots are non-adjacent, they have the same adjacencies to
$U$, and no sampled old-old edge has changed.  This gives the estimate.
\end{proof}

\begin{theorem}[Non-edge type $C_5$-free root-plantability]\label{thm:c5-nonedge-root}
The $C_5$-free graph class is root-plantable at the two-root non-edge type
$\eta$:
\[
  S_\eta=Q_\eta.
\]
\end{theorem}

\begin{proof}
Again $S_\eta\subseteq Q_\eta$ is the general support-passing statement.  Let
$\psi\in Q_\eta$, and let $U$ be a basic neighbourhood of $\psi$, specified by
$\eta$-flags $F_1,\ldots,F_s$ of size at most $m$ and a tolerance $\eps>0$.
Choose a convergent sequence of $C_5$-free graphs rooted at a non-edge,
\[
  (G_t,r_t,s_t),\qquad |G_t|\to\infty,
\]
with $p(\,\cdot\,,(G_t,r_t,s_t))\to\psi$.

Pick $\lambda>0$ so small that the constant in
\cref{lem:c5-nonedge-density-approx} satisfies $C_m\lambda<\eps/4$.  Put
\[
  L_t=\max(1,\lfloor\lambda |G_t|\rfloor),
  \qquad
  H_t:=P_{L_t}(G_t,r_t,s_t).
\]
By \cref{lem:c5-nonedge-planting-free}, every $H_t$ is $C_5$-free.  Passing to a
subsequence, assume the unlabelled graphs $H_t$ converge to some
$\phi_0\in Q_0$.

Let $R_t,S_t$ be the planted clusters.  A uniformly random ordered embedding of
the non-edge type into $H_t$ lands with its first label in $R_t$ and its second
label in $S_t$ with probability at least
\[
  \delta_t\ge
  \frac{|R_t||S_t|}{|V(H_t)|^2}
  \longrightarrow
  \left(\frac{\lambda}{1+2\lambda}\right)^2
  =:\delta>0,
\]
where the denominator is only an upper bound for the number of ordered non-edge
embeddings.  On this planted event, \cref{lem:c5-nonedge-density-approx} and the
convergence to $\psi$ imply, for all sufficiently large $t$,
\[
  |p(F_i,(H_t,x,y))-\psi(F_i)|\le\eps/2
  \qquad(1\le i\le s,\ (x,y)\in R_t\times S_t).
\]

Let
\[
  Y_\eta:=[0,1]^{\F^{\eta}[\T_{\mathrm{Graph}}]},
  \qquad
  \widetilde C=
  \{z\in Y_\eta:
    |z(F_i)-\psi(F_i)|\le\eps/2\text{ for all }i\}.
\]
Then $\widetilde C\cap X_\eta\subseteq U$, and the finite random non-edge-root
distribution of $H_t$ gives $\widetilde C$ mass at least $\delta/2$ for all
large $t$.  These finite random-root distributions converge weakly to
$\Ext_\eta(\phi_0)$ as measures on $Y_\eta$, so the closed-set form of the
Portmanteau theorem gives
\[
  \Ext_\eta(\phi_0)(\widetilde C)\ge\delta/2>0.
\]
Since $\Ext_\eta(\phi_0)$ is concentrated on $X_\eta$, its support intersects
$\widetilde C\cap X_\eta\subseteq U$.  As $U$ was arbitrary, $\psi\in S_\eta$.
\end{proof}

For comparison, ordinary blow-ups detect exactly the triangles.

\begin{lemma}\label{lem:c5-blowup}
Every independent blow-up of a $C_5$-free graph $G$ is again $C_5$-free if and only
if $G$ is triangle-free.
\end{lemma}

\begin{proof}
A $C_5$ in a blow-up projects to a closed walk of length $5$ in $G$, and conversely
any closed walk of length $5$ in $G$ lifts to a $C_5$ in a blow-up with classes of
size $\ge2$.  Since $G$ is $C_5$-free, such a walk---being of odd length $5$, hence
containing an odd cycle of length $3$ or $5$---exists if and only if $G$ has a
triangle.
\end{proof}

\subsection{The edge-type obstruction}

The same class is not root-plantable at the two-root edge type.  Let $\tau$ be
the type consisting of a labelled edge, and let $g\in\A^\tau[\T_{\mathrm{Graph}}]$
be the three-vertex $\tau$-flag in which the one unlabelled vertex is adjacent
to both labelled vertices; equivalently, $g$ is the triangle flag rooted at an
edge.

\begin{lemma}[Few triangles]\label{lem:c5-few-triangles}
If $G$ is $C_5$-free, then the number $T(G)$ of triangles in $G$ satisfies
\[
  T(G)\le\frac23 e(G).
\]
In particular, every unlabelled $C_5$-free limit has triangle density zero.
\end{lemma}

\begin{proof}
For every $v$, \cref{lem:c5-nbhd} gives $e(G[N(v)])\le d(v)$.  Summing over
vertices,
\[
  3T(G)=\sum_v e(G[N(v)])\le \sum_v d(v)=2e(G),
\]
since every triangle is counted once at each of its three vertices.  As
$e(G)=O(|V(G)|^2)$, the normalized triangle density is $O(1/|V(G)|)$ along any
sequence with $|V(G)|\to\infty$.
\end{proof}

\begin{corollary}[The edge-rooted triangle is pinned on ensembles]\label{cor:c5-edge-pinned}
Let $\phi_0\in Q_0$ be a $C_5$-free unlabelled homomorphism with
$\phi_0(\brak{\tau})>0$, and let $\phi^\tau\sim\Ext_\tau(\phi_0)$.  Then
\[
  \phi^\tau(g)=0
  \qquad\text{almost surely.}
\]
\end{corollary}

\begin{proof}
The downward average $\avg{g}{\tau}$ is a positive scalar multiple of the
unlabelled triangle.  By \cref{lem:c5-few-triangles},
$\phi_0(\avg{g}{\tau})=0$.  Razborov's extension identity gives
\[
  \E[\phi^\tau(g)]
  =
  \frac{\phi_0(\avg{g}{\tau})}{\phi_0(\brak{\tau})}
  =0.
\]
Since $\phi^\tau(g)\ge0$ pointwise, it is almost surely zero.
\end{proof}

\begin{definition}[Book graphs]\label{def:c5-book}
Let $B_n$ be the graph with two distinguished vertices $a,b$, the edge $ab$,
and $n$ further vertices each adjacent to both $a$ and $b$, with no edges among
the further vertices.
\end{definition}

\begin{lemma}[The book edge is $C_5$-free]\label{lem:c5-book}
Each $B_n$ is $C_5$-free.  Rooted at the ordered edge $(a,b)$, the sequence
$(B_n,a,b)$ has a convergent subsequence whose limit is a point $\psi\in Q_\tau$
with
\[
  \psi(g)=1.
\]
\end{lemma}

\begin{proof}
A leaf of $B_n$ has neighbourhood exactly $\{a,b\}$.  A simple cycle containing
a leaf must enter and leave it through $a$ and $b$.  Hence a simple cycle in
$B_n$ uses at most two leaves, so its length is at most $4$.  Thus $B_n$ has no
$C_5$.

For the rooted density, every unlabelled vertex of $B_n$ is adjacent to both
roots.  Therefore $p(g,(B_n,a,b))=1$ for all $n$, and any convergent subsequence
has the asserted limit.
\end{proof}

\begin{theorem}[The edge type is not root-plantable]\label{thm:c5-edge-not-root-plantable}
The $C_5$-free graph class is not root-plantable at the two-root edge type
$\tau$.  Concretely, $-g$ satisfies ensemble semantics but not quotient
semantics.
\end{theorem}

\begin{proof}
By \cref{cor:c5-edge-pinned}, every support point of every admissible
edge-rooted ensemble has $g=0$, so $-g$ is ensemble-nonnegative.  By
\cref{lem:c5-book}, there is a quotient homomorphism
$\psi\in Q_\tau$ with $\psi(g)=1$, hence $\psi(-g)=-1<0$.  Thus quotient
semantics fails for $-g$, and the support-closure criterion gives
$S_\tau\ne Q_\tau$.
\end{proof}

\begin{remark}
The $C_5$-free class therefore gives both a positive and a negative lesson.  A
single typical vertex in a positive-density $C_5$-free limit may have nontrivial
edge density to the rest of the graph, and \cref{thm:c5-one-root} shows that
every such rooted quotient limit can be planted.  A typical rooted edge is much
more rigid: the common-neighbour density is forced to vanish, while an
exceptional edge in a sparse book graph may have common-neighbour density one.
This is precisely a pinning obstruction, so the characterisation conjecture
survives this test but must be understood type by type.
\end{remark}

\section{Strengthening by a further constraint}

The framework is not tied to a single forbidden-subgraph constraint.  It applies
unchanged to any hereditary class, in particular to the intersection $\K\cap\K'$ of the
constrained class with a further graph property $\K'$, provided the intersection is again
hereditary.

\subsection{When the further constraint is hereditary}

If $\K'$ is itself a hereditary graph property---closed under induced subgraphs---then the
intersection $\T_1:=\K\cap\K'$ is hereditary, the quotient flag algebra $\A^\sigma[\T_1]$
is well-defined, and the support-closure criterion (\cref{thm:support-criterion}) applies
verbatim.  The equivalence of the two semantics is thus governed by root-plantability and
by nothing weaker: it does \emph{not} hold for all hereditary classes, as the degeneracy
obstruction (\cref{thm:degenerate-obstruction})---already in the $C_4$-free
instance---refutes equivalence under heredity alone.  The substantive task is therefore
to find checkable sufficient conditions for root-plantability.  Clone-closure is one
(\cref{thm:clone-root-plantable}); natural candidates for others are closure under
substitutions, true- or false-twin duplications, and amalgamation properties.

\subsection{When the further constraint is not hereditary}

If $\K'$ is not hereditary---for instance connectivity, or a parity condition on the
number of edges---then $\K\cap\K'$ is generally not closed under induced subgraphs.  This
breaks the basic starting point of Razborov's flag-algebra construction, which assumes
that every induced subgraph of a graph in the class is again in the class.  In this case
there may be no canonical quotient algebra $\A^\sigma[\K\cap\K']$ of the usual
flag-algebra kind.

One can still study constraints imposed on ambient homomorphisms, for example
conditions of the form
\[
  \phi(g_i)=0,
  \qquad
  \phi(h_j)\ge0.
\]
But this is better viewed as \emph{relative semantics inside the ambient flag
algebra} than as quotient semantics of a new hereditary class, and lies outside the
quotient/ensemble problem considered here.

\section{Open problems}

Three questions remain.

\begin{enumerate}[label=\textup{(\arabic*)}, leftmargin=2.8em]
  \item \textbf{Sufficient conditions beyond clone-closure.}
  The blow-up proof of \cref{thm:clone-root-plantable} duplicates labelled vertices into
  positive-mass clone classes.  A broader theorem may permit more general substitutions
  or local replacements, as long as the labelled finite graph can be amplified into
  positive measure without leaving the constrained class; closure under true- or
  false-twin duplication is a natural first generalisation.

  \item \textbf{The characterisation conjecture.}
  \Cref{conj:characterisation} asks whether the absence of a pinning obstruction is
  already sufficient for root-plantability.  The $C_5$-free split sharpens the test:
  the one-vertex type has no elementary pinning obstruction and is root-plantable, while
  the edge type fails for exactly a pinning reason.  A decisive next example would be a
  type or class with no pinning obstruction but still no root-plantability, or a theorem
  ruling this out under checkable hypotheses.

  \item \textbf{Whether the gap ever reaches a density bound.}
  By \cref{thm:no-interior} every pinning obstruction is to a boundary value, and
  \cref{sec:empty-type} shows the exhibited obstructions inert for density bounds.  A gap
  that did affect a genuine density bound would require a non-root-plantable type whose
  labelled support is rich enough to carry a non-negative element that is no restricted sum
  of squares.  None is known; the dense $C_5$-free labelled limits, which do not collapse to a
  point, are again the natural place to look.
\end{enumerate}

\begin{thebibliography}{9}

\bibitem{Razborov2007}
A.~A.~Razborov.
\newblock Flag algebras.
\newblock \emph{Journal of Symbolic Logic}, 72(4):1239--1282, 2007.

\bibitem{KovariSosTuran}
T.~K\H{o}v\'ari, V.~T.~S\'os, and P.~Tur\'an.
\newblock On a problem of K. Zarankiewicz.
\newblock \emph{Colloquium Mathematicae}, 3:50--57, 1954.

\bibitem{BondySimonovits}
J.~A.~Bondy and M.~Simonovits.
\newblock Cycles of even length in graphs.
\newblock \emph{Journal of Combinatorial Theory, Series B}, 16(2):97--105, 1974.

\end{thebibliography}

\end{document}
